\documentclass{article}
\usepackage[T1]{fontenc}
\usepackage{lmodern}
\usepackage{graphicx}
\usepackage{amsmath,amssymb}
\usepackage[a4paper,margin=2cm]{geometry}
\usepackage[absolute,overlay]{textpos}
\setlength{\TPHorizModule}{1mm}
\setlength{\TPVertModule}{1mm}
\usepackage[indonesian]{babel}
\usepackage{hyperref}
\setlength{\emergencystretch}{2em}
% unit: Halaman 4

\begin{document}

% === Halaman 4 (python/native) ===

• Oleh karena itu, dikenal dua macam cipher di dalam kriptografi klasik:

1. Cipher Substitusi (substitution Cipher)

       - metode enkripsi dan dekripsi menggunakan teknik substitusi

2.  Cipher Transposisi (transposition Cipher)

        - metode enkripsi dan dekripsi menggunakan teknik transposisi

• Kombinasi kedua teknik tersebut membentuk product cipher  atau 

super enkripsi

          product cipher = cipher substitusi + cipher transposisi

4

\end{document}
